% language=us signal=signal

% \nologbuffering

\usemodule[present-boring,abbreviations-logos]

\definelayer
  [extratext]
  [preset=rightbottom,
   height=\textheight,
   width=\textwidth]

\setupbackgrounds
  [text]
  [background=extratext]

\startdocument
  [title={blowing the whistle},
   banner={tagged or flagged},
   location={bachotex\enspace {\bf 2025}\enspace meeting}]

\starttitle[title=Outline]

\startitemize
\startitem
    Tagged \PDF\ gets enforced upon us. Validators dominate over properly designed
    documents. What happened to rendering for specific target groups and
    purposes?
\stopitem
\startitem
    Universities want documents to comply so they offer checkers. We managed to
    get them happy (after all, they are clueless about content) except for
    colors; show an example.
\stopitem
\startitem
    However, the tool that we use is rather inaccessible. Take error messages.
    How accessible are they?
\stopitem
\startitem
    But \unknown\ let's test this! We introduce several techniques, for instance
    QR codes. After all, this is the time of phones and we need to attract youth.
    Of course it now can become fashion to generate errors.
\stopitem
\startitem
    But isn't it more challenging to use colors? How can they help us? We need to
    do some research!
\stopitem
\startitem
    Tove's color \PDF\ talk. A scientific experiment to determine if colors in
    \PDF\ are really an issue. If not, we can move back to \TEX.
\stopitem
\startitem
    Back to \TEX. We're going to do two tests: with QR tags and with colors. And
    in the end the user might decide. We also explore sound as feedback!
\stopitem
\stopitemize

\stoptitle

\stopdocument
